\section{Introduction}

NEMA targets a narrow but recurrent problem in hardware-oriented compilation: how to specify, implement, and validate neural-graph computations such that (i) executions are deterministic by construction and (ii) equivalence across implementations can be checked mechanically. In this project, the computational model is a sparse directed graph of neurons and synapses evaluated in discrete ticks. Without an explicit contract, two implementations can diverge silently due to update-order effects, floating-point drift, or ``almost the same'' numerics inside different toolchains. NEMA's design centers on making these failure modes observable and rejectable through a small set of gates that operate on stable, machine-readable artifacts, in line with broader observations on reproducibility variance in modern ML systems \cite{PhamASE2020Variance,ChenICSE2022ReproDL,ZhuangArxiv2106.11872,ShanmugaveluArxiv2408.05148}.

At the semantic level, NEMA v0.1 defines a deterministic tick semantics (\texttt{nema.tick.v0.1}) with a snapshot rule that guarantees order-independent results with respect to update ordering, while still evaluating nodes in a well-defined index order. The numeric model is fixed-point: saturating overflow is the only permitted overflow behavior, and rounding is restricted to round-to-nearest ties-to-even (RNE). Nonlinearities are mediated by a fixed tanh lookup-table policy (\texttt{nema.tanh\_lut.v0.1}). These choices reduce the semantic surface area of implementation-defined behavior and allow NEMA to define bit-exactness in terms of a canonical digest.

Bit-exactness is defined over per-tick state digests. For each tick, NEMA collects node voltages \texttt{V} in node-index order, packs each raw voltage as a signed int16 little-endian value (Q8.8), and computes SHA-256 over the packed byte array. Two executions are bit-exact equal when all per-tick digests match. This definition is independent of any particular simulator backend and yields a simple, stable contract that can be applied uniformly to software runners and hardware-oriented flows.

NEMA programs are represented in a validated IR (JSON shape mirrored by \texttt{nema\_ir.proto}), with an explicit validator enforcing structural invariants (e.g., unique node/edge IDs, edge references to existing nodes, chemical edges directed, gap edges symmetric, non-negative conductances) and integrity constraints for external references. In particular, IR graphs may be inlined or referenced via \texttt{graph.external}; referenced bundles must exist on disk, and if a \texttt{sha256} is provided (not a placeholder token) it must match the file digest. This makes ``what graph did you actually run?'' part of the executable contract rather than an implicit assumption.

The repository provides an end-to-end toolchain around this contract: a textual DSL front-end that lowers deterministically to the IR, a reference simulator (\texttt{nema/sim.py}) and fixed-point library (\texttt{nema/fixed.py}), a hardware test harness (\texttt{nema/hwtest.py}), an HLS code generator (\texttt{nema/codegen/hls\_gen.py}), and a report schema (\texttt{tools/bench\_report\_schema.json}). The primary entrypoint for reproducibility is the self-audit gate \texttt{tools/audit\_min.py}, which evaluates software and hardware modes and emits a JSON decision together with criteria that must hold (e.g., digest stability, bench verification success, toolchain availability, and presence of hardware report evidence).

Our minimal evaluation is defined as a set of benchmark manifests and machine-checkable metrics. The paper core includes B1 (a micro-graph sanity benchmark), B3 (a connectome-sized kernel with targetId ``CE/302-7500''), and B4 (an external connectome-bundle workflow with targetId ``CE/8-12''). Correctness is evaluated by \texttt{nema bench verify}, which reports \texttt{ok=true} and \texttt{mismatches=[]} (mismatch=0) for B1 and B3 on the traces evaluated by their manifests. Hardware evidence is evaluated through \texttt{audit\_min --mode hardware}, which must return \texttt{decision=\"GO\"} and, by construction of the gate, implies toolchain detection, successful execution of the hardware pipeline (G0b), capture of QoR/report evidence (G2), and presence of Vivado implementation timing with non-null WNS (G3) in the generated \texttt{bench\_report.json} artifacts. Table-style summaries of these gate outcomes are included in the artifact pack (e.g., \texttt{papers/paperA/artifacts/tables/gates\_summary.*}).

Finally, NEMA's artifact story is intentionally operational: the paper's claims correspond to scripts that produce stable outputs in stable paths. The artifact pack records toolchain versions (Vivado/Vitis HLS), hashes the artifact manifest, and provides precomputed gate evidence (software and hardware audit JSON). The project's current scope explicitly excludes claims of ML accuracy advantages, biological fidelity, or on-board power/latency measurements; those are out of scope for the verified contract presented here.

\subsection{Contributions}
\begin{itemize}
\item \textbf{Deterministic semantics + numeric contract}: We specify a deterministic tick semantics (snapshot rule + index-ordered iteration) and a fixed-point numeric model (SATURATE overflow, RNE rounding), together with a bit-exact definition based on per-tick SHA-256 digests of packed Q8.8 voltages. (Artifact: \texttt{spec.md} extracts; Artifact: \texttt{nema\_ir.proto})
\item \textbf{Script-verifiable bit-exactness on core benches}: We provide a manifest-driven verification command (\texttt{nema bench verify}) that reports \texttt{mismatches=[]} (mismatch=0) for B1 and B3 on the evaluated traces, making cross-implementation equivalence machine-checkable. (Artifact: \texttt{build/paper\_verify/b1/*} and \texttt{build/paper\_verify/b3/*} produced by the commands)
\item \textbf{Hardware evidence captured and gated}: We integrate HLS execution + QoR capture (G0b/G2) and Vivado implementation timing capture with non-null WNS (G3) into the same reporting pipeline, and expose a single gate (\texttt{audit\_min --mode hardware}) that must return \texttt{decision=GO}. (Table: \texttt{papers/paperA/artifacts/tables/gates\_summary.md}; Artifact: \texttt{papers/paperA/artifacts/evidence/audit\_hardware.json}; Artifact: \texttt{build\_hw/**/bench\_report.json})
\item \textbf{External bundle workflow as a first-class contract}: We support IR graphs referenced via \texttt{graph.external} with file-existence and optional SHA-256 verification, and include a benchmark (B4) that exercises this external connectome-bundle path through the same manifest/report interface. (Artifact: \texttt{benches/B4\_real\_connectome/manifest.json}; Artifact: \texttt{build/bench\_verify\_eqlozbf7/B4\_celegans\_external\_bundle/bench\_report.json})
\item \textbf{Reproducible artifact packaging with hashes and versions}: We package gate evidence and provenance metadata (git hash, toolchain versions, SHA-256 of \texttt{artifact\_manifest.json}, and counts of tables/figures/evidence) to make the verification story auditable from files alone or by re-running scripts. (Artifact: \texttt{papers/paperA/artifacts/artifact\_manifest.json} + recorded SHA-256; Table: \texttt{gates\_summary.*}; Fig: \texttt{papers/paperA/artifacts/figures/gates\_status.txt})
\end{itemize}
